\section{\bt{Publication, Maintenance, and Optimization}{发布、维护与优化}}
\label{sec:maintenance}
\subsection{\bt{Publication checks on successful sessions}{成功会话的发布前检查}}
\bi{Candidates are extracted only from tasks whose question states the business definition and whose answer matched the reference answer. The extractor sees the question, the SQL and arithmetic behind the answer, rejected queries, and metadata, but no business rows or result values. Seven deterministic publication checks follow: a stated business definition; a match with the reference answer, with a recorded calculation that reproduces it from the query results; static checks (columns exist, the measure uses only permitted columns, joins follow validated paths, and required grain filters are present without question literals); key uniqueness on the current snapshot; the SQL check (no join known to fan out, no missing required filter); replay of the example SQL; and replay of the generated SQL compiled from the fields. The replay of the generated SQL checks what other agents will actually read, so a candidate that omits a filter or hard-codes a parameter fails it. Learning questions supply the publication record; held-out questions only measure reuse.}
{只从题目给出业务定义且答案与参考答案一致的任务中提取候选。提取器能看到题面、答案背后的 SQL 与算式、被拦下的查询以及元数据，但看不到业务行和结果值。随后执行七项确定性的发布前检查：给出了业务定义；答案与参考答案一致，且记录的算式能从查询结果复现该答案；静态检查（列存在、度量只用允许的列、连接沿已验证路径、包含必需的粒度过滤且不含题目字面量）；当前快照上的键唯一；SQL 检查（没有已知会放大行数的连接，不缺必需的过滤）；示例 SQL 重放；由字段编译的生成 SQL 重放。生成 SQL 的重放检查的是其他智能体实际会读到的内容，因此遗漏过滤或把参数写死的候选会失败。学习题提供发布记录，留出题只用于测量复用效果。}

\subsection{\bt{Selective recheck and reuse}{选择性重查与复用}}
\bi{On use, maintenance compares the versions recorded for the tables each condition reads with the current versions. Conditions over unchanged tables keep their check results. For each remaining condition, \system\ looks up a completed check result for the same normalized form at the current versions; otherwise it runs $Q_c$, and concurrent requests for the same check share that one execution (request coalescing). Joins are checked through the shared join entries, so a changed or vanished join path also makes the definition pending.}
{使用时，维护过程把每个条件所读各表记录下的版本与当前版本比较。未变化表上的条件保留原有检查结果。对其余条件，\system\ 查找同一规范化形式在当前版本上的已完成检查结果；否则执行 $Q_c$，对同一检查的并发请求共享这一次执行（请求合并）。连接通过共享的连接条目检查，因此连接路径被改变或消失时，定义同样变为待验证。}

\subsection{\bt{Invalidation and automatic repair}{失效与自动修复}}
\bi{A failed condition invalidates the revision and records the reason. Repair handles two cases: a failed fact-table grain, for which it searches equality filters on low-cardinality columns, and a changed join path, whose new filters it adopts. A grain filter must keep exactly one row per business key without losing keys. The fix must also be \emph{unique}: \system\ enumerates every passing filter and repairs only if exactly one exists, because two complete copies of the data, such as a primary table and an appended backup, each pass the structural tests, and nothing structural tells which holds the business truth. A unique fix must then pass the static, grain, and SQL checks and a regression test that runs the replacement and the agent's learning SQL as \emph{as-of queries} (time-travel queries) at learning time and requires equal results: the agent's SQL was verified correct only on that data, and on current data it can itself be stale, counting old versions after a restatement. Warehouses answer as-of queries through time travel; our prototype keeps a copy of the tables taken at publication. Repair searches over the same table and key are shared across definitions, and a failed repair leaves the definition unavailable until it is relearned.}
{条件失败会使版本失效并记录原因。修复只处理两种情况：事实表粒度失效时，搜索低基数列上的等值过滤；连接路径改变时，采用新路径的过滤。粒度过滤必须使每个业务键恰好保留一行且不丢键。修复还必须是\emph{唯一}的：\system\ 列出全部通过检查的过滤，只有恰好一个时才修复，因为主表与追加的备份副本这样的两份完整数据各自都能通过结构检查，结构上无从判断哪一份承载了业务事实。唯一的修复还要通过静态、粒度与 SQL 检查以及回归测试：把替代版本与智能体的学习 SQL 一起以学习时刻为基准作 \emph{as-of 查询}（时间旅行查询），要求结果相同；智能体的 SQL 只在那时的数据上被判过对，在当前数据上它本身可能已经过期，例如更正保留旧行后会把旧行也算进去。数仓用 time travel 回答 as-of 查询；我们的原型在发布时保留一份表的副本。同一张表、同一个键上的修复搜索在定义之间共享；修复失败的定义保持不可用，直到重新学习。}

\begin{algorithm}[t]
\caption{\bt{Automatic repair of a failed grain condition}{粒度条件失效后的自动修复}}
\label{alg:repair}
\footnotesize
\begin{algorithmic}[1]
\Require \bt{revision $m^r$ whose $\operatorname{KeyUnique}(F,K,\varphi)$ failed; learning question $x$; the agent's learning SQL $R$; the database $D_\ell$ as of learning time; low-cardinality columns $A$ of $F$}{$\operatorname{KeyUnique}(F,K,\varphi)$ 失效的版本 $m^r$；学习题 $x$；智能体的学习 SQL $R$；学习时刻的数据库 $D_\ell$（as-of）；$F$ 的低基数列 $A$}
\State $P\gets\emptyset$
\ForAll{$a\in A$ \textbf{and} $v\in\pi_a(F)$}
  \State $\psi\gets\varphi\wedge(a=v)$
  \If{$\operatorname{KeyUnique}(F,K,\psi)$ \textbf{and} $\pi_K(\sigma_\psi F)=\pi_K(\sigma_\varphi F)$}\label{alg:repair:keys}
    \State $P\gets P\cup\{\psi\}$ \Comment{\bt{one row per key, no key lost}{每键一行且不丢键}}
  \EndIf
\EndFor
\If{$|P|\neq 1$}\label{alg:repair:unique}
  \State \Return \textsc{invalidate} \Comment{\bt{no fix, or several equivalent fixes}{没有修复，或多个等价修复}}
\EndIf
\State $m'\gets m^r$ \bt{with}{，令} $\varphi:=\psi$ \bt{for the single $\psi\in P$; re-derive $C(m')$}{（$\psi$ 为 $P$ 的唯一元素），重新推导 $C(m')$}
\If{$m'$ \bt{fails the static, grain, or SQL checks}{未通过静态、粒度或 SQL 检查} \textbf{or} $\operatorname{compile}(m',x)(D_\ell)\neq R(D_\ell)$}\label{alg:repair:regression}
  \State \Return \textsc{invalidate} \Comment{\bt{regression test}{回归测试}}
\EndIf
\State \bt{publish}{发布} $m^{r+1}=m'$
\end{algorithmic}
\end{algorithm}


\bi{Algorithm~\ref{alg:repair} summarizes the grain case: a published replacement restores one row per business key without losing a key on the repair snapshot (line~\ref{alg:repair:keys}), is the only filter in the search space that does so (line~\ref{alg:repair:unique}), and reproduces the agent's learning answer as of learning time (line~\ref{alg:repair:regression}). Uniqueness refuses structurally equivalent populations even when the reference would accept the wrong one; regression refuses a filter that changes the learning answer on the data the agent learned from.}
{算法~\ref{alg:repair} 概括了粒度失效的情形：被发布的替代版本在修复所在快照上恢复每个业务键恰好一行且不丢键（第~\ref{alg:repair:keys} 行），是搜索空间中唯一做到这一点的过滤（第~\ref{alg:repair:unique} 行），并按学习时刻复现智能体的学习答案（第~\ref{alg:repair:regression} 行）。唯一性拒绝结构上等价的多个总体，哪怕参照会接受其中错误的一个；回归拒绝在智能体学习所用的数据上改变学习答案的过滤。}

\bi{Requests that arrive while a repair runs join the running repair instead of being rejected, so a successful repair serves them as well.}
{修复期间到达的请求不被拒绝，而是与进行中的修复合并，因此修复成功时它们同样得到服务。}

\subsection{\bt{Optimization revisions}{优化版本}}
\label{sec:optimization}
\bi{A published definition is correct but not necessarily cheap, and because agents copy its generated SQL, a cheaper equivalent benefits every consumer. After learning, \system\ runs an optimization pass over the valid definitions; candidates are structured definitions, so compilation, the conditions, and Proposition~\ref{prop:sufficiency} apply to them. Two rules propose candidates: \emph{key range} replaces the period join with a range predicate on the fact table's date key, whose bounds scalar subqueries read from the date dimension, and \emph{left-join removal} drops a left join that contributes no column and no filter. Only when no rule candidate is published does an LLM propose candidates from the definition, its generated SQL and plan, and table sizes. A candidate must pass the static, grain, and SQL checks and three tests. \emph{Assumption}: the key-range rule adds \emph{date-key contiguity}, which requires the keys of every month to form an integer interval that begins one above the previous month's last key; it is checked before adoption and maintained afterwards. \emph{Equivalence}: candidate and current revision return equal results for the learning question as of learning time and, within one repeatable-read snapshot, for the twelve months of the learning year, a quarter, a period difference, and a monthly ranking. \emph{Cost}: in the same snapshot, the paired execution times of at least eight periods, measured in alternating order, must have a 95\% confidence interval below zero and a mean saving of at least 10\%. For the rules, equivalence does not rest on the samples.}
{已发布的定义是正确的，但不一定是最省的；由于智能体会照抄它的生成的 SQL，一个更省的等价写法对每个使用者都有收益。学习结束后，\system\ 对有效的定义执行一轮优化；候选都是结构化定义，因此编译、条件和命题~\ref{prop:sufficiency} 都适用于它们。两条规则提出候选：\emph{日期键范围}把期间连接改为事实表日期键上的范围谓词，上下界由标量子查询从日期维度读出；\emph{去掉左连接}去掉一个既不提供列也不带过滤的左连接。只有在没有规则候选发布时，才由大模型根据定义、其生成的 SQL 与执行计划以及表的规模提出候选。候选必须通过静态、粒度和 SQL 检查以及三项测试。\emph{前提}：日期键范围规则新增\emph{日期键连续}条件，要求每个月的键构成一个整数区间，且起点比上个月的最后一个键大一；采纳前检查，之后受维护。\emph{等价}：候选与当前版本对学习题按学习时刻返回相同结果，并在一个可重复读快照内对学习年份的 12 个月、一个季度、一次期间差值和一次按月排名返回相同结果。\emph{代价}：在同一快照中，至少 8 个期间以交替顺序测得的配对执行时间差，其 95\% 置信区间必须整体低于 0，平均节省至少 10\%。对规则而言，等价并不依赖抽样。}
\begin{proposition}[\bt{Rule equivalence}{规则等价}]
\label{prop:rewrite}
\bi{Let the date dimension's key be unique on $D$. (a) If date-key contiguity holds on $D$, then for every period consisting of consecutive months, the generated SQL of a key-range revision returns on $D$ what the generated SQL of its previous revision returns. (b) If the right side of a removed left join is unique on the join key on $D$, the generated SQL without the join returns on $D$ what the generated SQL with it returns.}
{设日期维度的键在 $D$ 上唯一。（a）若日期键连续在 $D$ 上成立，则对每个由连续月份组成的期间，日期键范围版本的生成的 SQL 在 $D$ 上返回其前一版本的生成的 SQL 所返回的值。（b）若被去掉的左连接的右侧在连接键上于 $D$ 上唯一，则去掉该连接后的生成的 SQL 在 $D$ 上返回原生成的 SQL 所返回的值。}
\end{proposition}
\begin{proof}[\bt{Proof sketch}{证明概要}]
\bi{(a) By contiguity, the keys of the months of a period form exactly the integer interval between the bounds the subqueries read, so a fact row's date key lies in the interval if and only if it joins a date row of the period, and by uniqueness only one; a key that matches no date row lies outside the dimension's key range, and both forms drop it. (b) A left join with at most one match keeps each left row exactly once and supplies nothing that the aggregate or the filters read.}
{（a）由连续性，一个期间所含月份的键恰好构成子查询读出的上下界之间的整数区间，因此事实行的日期键落在区间内当且仅当它连接到该期间的一行日期记录，且由唯一性只有一行；匹配不到任何日期行的键落在维度的键范围之外，两种写法都排除它。（b）至多匹配一行的左连接把每个左表行恰好保留一次，也不提供聚合或过滤读取的任何内容。}
\end{proof}
\bi{The assumption of (a) is a condition of the new revision; that of (b) is a condition of the previous revision that the new revision no longer needs, since if the right side later loses uniqueness, the previous revision becomes invalid while the new revision stays correct by Proposition~\ref{prop:sufficiency}. Wherever the maintained conditions hold, a rule revision thus agrees with its previous revision on every period; the sampled tests check the implementation and measure the cost, and only an LLM-proposed rewrite rests on them for equivalence. Maintenance checks an added assumption last; when it alone fails, every condition shared with the previous revision has just been confirmed, and \system\ invalidates the optimized revision and restores the previous revision or, when a repair has since replaced the optimized revision and no previous revision is kept, undoes the rule (Sections~\ref{sec:opt-eval} and~\ref{sec:lifecycle-eval}).}
{（a）的前提是新版本的条件；（b）的前提是前一版本的条件，新版本不再需要它：若右侧之后失去唯一性，前一版本失效，而新版本由命题~\ref{prop:sufficiency} 仍然正确。因此，只要受维护的条件成立，规则版本就在每个期间上都与前一版本一致；抽样测试检查的是实现并测量代价，只有大模型提出的改写依赖它们来确立等价。维护最后检查新增的前提；只有它不成立时，与前一版本共有的每个条件都刚刚得到确认，\system\ 使优化后的版本失效，恢复前一版本；若优化后的版本此后经过修复、已不保留前一版本，则撤销改写规则（第~\ref{sec:opt-eval} 与~\ref{sec:lifecycle-eval} 节）。}
